\section{Experiments}
\label{sec:method:experiment}

\subsection{Research Design}
\label{sec:method:experiment:design}
The manipulated variable is the reduction method, with encoder, dataset, splits, label,
optimizer and evaluation protocol held constant against the full-graph model of RQ1 as
reference. RQ1 needs the reference point trainable, bounding corpus scale
(\ref{sec:method:experiment:scale}). RQ2 needs one batching budget
(\ref{sec:method:experiment:batching}). RQ3 needs hyperparameters fixed at the full-graph
values (\ref{sec:method:experiment:hp}). RQ4 needs pairs matched on \emph{achieved}
compression. RQ5 needs only a second evaluation pass over the Phase~3 checkpoints. H1 (path
contraction) and H2 (heuristic strength) are stated with their questions
(\ref{sec:intro:rq3},~\ref{sec:intro:rq4}).

% Original, condensed above:
% The manipulated variable is the reduction method applied to the training data; the encoder, dataset, splits, label,
% optimizer and evaluation protocol are held constant, so that differences in outcome are
% attributable to the reduction. The reference point is the full-graph model of RQ1, and every
% reduced configuration is reported relative to it rather than in isolation.
%
% Each question makes its own demand on that design. RQ1 requires the unreduced reference
% point to be trainable at all, which is what bounds the corpus scale
% (\ref{sec:method:experiment:scale}). RQ2 requires every configuration to be measured under
% one batching budget, since memory and throughput are not comparable across budgets
% (\ref{sec:method:experiment:batching}). RQ3 requires the reduction to be the only thing that
% changed, which is why hyperparameters are fixed at the full-graph values
% (\ref{sec:method:experiment:hp}). RQ4 requires pairs matched on \emph{achieved} compression
% rather than on configured parameters, so that a retention gap is not a compression gap in
% disguise. RQ5 requires no new training, only a second evaluation pass over the Phase~3
% checkpoints. The two hypotheses under test, H1 on path contraction and H2 on the strength of
% the heuristics, are stated with the questions they belong to
% (\ref{sec:intro:rq3},~\ref{sec:intro:rq4}).

\paragraph{Positive control.}
Colour refinement at count-cap $\infty$ merges only nodes the encoder provably cannot
distinguish (\ref{sec:method:architecture:exact}), so a model trained on it must score
identically to the full-graph baseline on both passes: any deviation is a defect, not a
finding, and without this control RQ5 cannot separate a generalization failure from an
evaluation-path bug.

% Original, condensed above:
% One configuration has a predetermined outcome. Colour refinement at count-cap $\infty$ on
% the exact-compression track merges only nodes the encoder provably cannot distinguish
% (\ref{sec:method:architecture:exact}), so a model trained on those coarsened graphs must
% score identically to the full-graph baseline, on reduced and on unreduced graphs alike. Any
% deviation is a defect to be found rather than a finding to be reported, and without it RQ5
% cannot separate ``models trained on reduced graphs do not generalize'' from ``the evaluation
% path is wrong''.

\paragraph{The experiments.}
\label{sec:method:experiment:phases}
Four experiments implement the design and do not map one-to-one onto the five research
questions (\ref{tab:experiment_phases}). RQ4 re-analyses the Phase~3 runs as a paired
comparison rather than a separate experiment, costing no additional training, and RQ5 reuses
the Phase~3 checkpoints for only an evaluation pass. The ordering is a dependency chain:
Phase~3 needs Phase~2's cached reduction artifacts, and Phase~4 needs Phase~3's checkpoints.

% Original, condensed above:
% Four experiments implement the design, and they do not map one-to-one onto the five research
% questions (\ref{tab:experiment_phases}). RQ4 is a re-analysis of the Phase~3 runs under a
% paired comparison rather than a separate experiment, which is why it costs no additional
% training; RQ5 reuses the Phase~3 checkpoints and adds only an evaluation pass. The ordering
% is a dependency chain rather than a preference: Phase~3 cannot start until Phase~2 has
% produced the cached reduction artifacts, and Phase~4 needs Phase~3's checkpoints.

\begin{table}[htbp]
\centering
\caption{The four experimental phases, what each answers, and what each costs. Only
Phases~1 and~3 consume training compute.}
\label{tab:experiment_phases}
\small
\begin{tabular}{llll}
\toprule
\textbf{Phase} & \textbf{Answers} & \textbf{What is run} & \textbf{Training cost} \\
\midrule
1 & RQ1      & Full-graph model and baselines        & one run per model \\
2 & RQ2      & Offline reduction precompute, profiled & none \\
3 & RQ3, RQ4 & One model per reduced configuration    & one run per configuration \\
4 & RQ5      & Phase~3 checkpoints on full graphs     & none (evaluation only) \\
\bottomrule
\end{tabular}
\end{table}

Phase~1 also sets the baseline hardware cost (peak memory, epoch time) every reduced
configuration is measured against. Phase~3 places reduced configurations on a Pareto front
of predictive degradation versus memory and speed. Phase~4 tests structural generalization
on unreduced graphs. Encoder family, target synthesis algorithm and graph scale are held
fixed throughout, so the design ranks reduction methods under these conditions only
(\ref{sec:discussion:limitations}).

% Original, condensed above:
% Phase~1 establishes both the predictive reference point and the baseline hardware cost
% (peak memory, epoch time) that every reduced configuration is measured against. Phase~2
% requires no training and determines which methods are viable at corpus scale at all.
% Phase~3 places the reduced configurations against the Phase~1 reference on a Pareto front
% of predictive degradation versus memory and speed. Phase~4 queries models trained
% exclusively on reduced graphs with unreduced ones, testing structural generalization.

\subsection{Setup}
\label{sec:method:experiment:setup}

\subsubsection{Data Splitting}
\label{sec:method:experiment:splitting}
An 80/10/10 split by volume, partitioned \emph{by design}: all graphs from a given source
circuit fall in the same split, since a random split would place a base graph in training
and its near-identical optimized descendant in test, letting the model score by recognising
the design rather than reading its structure. The cost is a harder task: the test set
contains circuits never seen in any state. Two alternatives, splitting by recipe and a
plain random split, are implemented but not used for the reported results, run once each as
the protocol-sensitivity experiment RQ1a (\ref{sec:results:rq1:protocol}), each an
established evaluation setting in the literature (Table~\ref{tab:split_protocols}).

% Original, condensed above:
% An 80/10/10 split by volume, partitioned \emph{by design}: all graphs derived from a given
% source circuit fall in the same split. This is the strong choice and it is deliberate.
% Because the corpus is built by transforming and then repeatedly optimizing a small number of
% source designs, a random split would place a base graph in training and its own optimized
% descendant in test: structurally near-identical circuits in different optimization states.
% The model could then score well by recognising the design rather than by reading its
% structure. Splitting by design removes that path entirely, at the cost of making the task
% harder: the test set contains circuits the model has never seen in any state.
%
% Two alternative strategies are implemented and available (split by optimization recipe, and
% a plain random split). They are not used for the reported results; they are run once, at the
% primary configuration, as the protocol-sensitivity experiment RQ1a
% (\ref{sec:results:rq1:protocol}).
% The three protocols differ in what they hold out, and each corresponds to an evaluation
% setting already established in the literature (Table~\ref{tab:split_protocols}).

\begin{table}[htbp]
\centering
\caption{The three train/test protocols implemented in this work, leakiest to strictest.
Only the design-disjoint protocol is used for the reported results. The other two quantify
how much of a score is design recognition rather than structural reading.}
\label{tab:split_protocols}
\small
\begin{tabular}{@{}p{0.15\linewidth}p{0.30\linewidth}p{0.27\linewidth}p{0.12\linewidth}@{}}
\toprule
\textbf{Protocol} & \textbf{Held out} & \textbf{Established as} & \textbf{Role here} \\
\midrule
Random (per-row)
& Nothing: graphs are assigned independently, so a base graph and its descendant may
  straddle the split.
& Common default in circuit representation-learning evaluations.
& Ablation \\[4pt]
Recipe-disjoint
& Whole recipes, across all designs: circuits stay seen, the recipe does not.
& OpenABC-D Variant~1 \citep{chowdhury2021openabc}; LOSTIN transductive
  \citep{wu2022lostin}.
& RQ1a \\[4pt]
Design-disjoint
& Whole source circuits: every graph from a held-out design, in any optimization state.
& OpenABC-D Variant~2 \citep{chowdhury2021openabc}; LOSTIN inductive
  \citep{wu2022lostin}.
& \textbf{Headline} \\
\bottomrule
\end{tabular}
\end{table}

% TODO: the random-split and recipe-split rows are placeholders until those runs land.

\subsubsection{Graph Scale Bounds}
\label{sec:method:experiment:scale}
The corpus is bounded at 366,040 nodes per AIG\@, the scale that makes RQ1 possible: a
single unreduced graph fits in accelerator memory, so the full-graph baseline can actually
be trained. Larger circuits would remove that reference point, leaving only reduced
configurations to compare with each other. This is a scoping decision with a cost: the
motivation invokes circuits of millions to billions of nodes, which the evaluation does not
reach (\ref{sec:discussion:limitations:scalability}).

% Original, condensed above:
% The corpus is bounded at 366,040 nodes per AIG\@. This intermediate scale is what makes RQ1
% possible: a single unreduced graph fits in accelerator memory, so the full-graph baseline
% every other result is measured against can actually be trained. Larger circuits would remove
% the reference point and leave only reduced configurations to compare with each other.
%
% The bound is a scoping decision with a cost, and it is the limitation closest to the
% motivation: the argument for this work invokes circuits of millions to billions of nodes, and
% the evaluation does not reach them
% (\ref{sec:discussion:limitations:scalability}).

\subsubsection{Dynamic Batching}
\label{sec:method:experiment:batching}
Graphs vary in size by orders of magnitude, so a fixed batch size in graphs gives wildly
varying memory use and fails on the largest circuits. Batches are instead assembled to a
maximum \emph{node} budget by greedy bin-packing, bounding activation memory regardless of
which graphs are drawn. The plan is computed once and cached. Two deliberately distinct
budgets are used, training the smaller and evaluation (no activations or gradients stored) a
substantially larger one, kept separate to avoid invalidating the batch plan the trained
checkpoints were produced under.

\paragraph{Constraint on efficiency measurements.}
The evaluation budget must be identical across all configurations, since peak memory and
throughput both depend on it: a mixed set of budgets would make the RQ2 comparison
meaningless, not merely noisy, constraining how~\ref{sec:results:rq2} may be read. One
irreducible peak remains regardless: a graph larger than the budget still forms a batch of
its own, since a graph-level target cannot be split, so the largest circuit sets a floor on
peak memory for any configuration that does not reduce node counts.

% Original, condensed above:
% Graphs in this corpus vary in size by orders of magnitude, so a fixed batch size in graphs
% gives wildly varying memory use and eventually fails on the largest circuits. Batches are
% instead assembled to a maximum \emph{node} budget by greedy bin-packing, which bounds
% activation memory regardless of which graphs are drawn. The batch plan is computed once and
% cached, since replanning over a corpus this size is expensive.
%
% Two budgets are used and they are deliberately distinct. Training uses the smaller of the
% two; evaluation, which stores no activations or gradients, uses a substantially larger one.
% Keeping them separate rather than raising a single shared value avoids invalidating the batch
% plan the trained checkpoints were produced under.
%
% Constraint on efficiency measurements: the evaluation budget must be identical across all
% configurations, because peak memory and throughput both depend on it. Any change requires
% re-running every configuration; a mixed set of budgets makes the RQ2 comparison meaningless
% rather than merely noisy. This is stated here because it constrains how the results
% in~\ref{sec:results:rq2} may be read and combined.
%
% One irreducible peak remains that no budget can lower: a graph larger than the budget still
% forms a batch of its own, since a graph-level target cannot be split across batches. The
% largest circuit in the corpus therefore sets a floor on peak memory for every configuration
% that does not reduce node counts.

\subsubsection{Model Configuration}
\label{sec:method:experiment:model}
An Edge-Aware Graph Convolutional Network (GCN+) predicts optimizability under
\texttt{Orchestrate}, described in~\ref{sec:method:architecture} and held fixed across all
reduction configurations. Baselines are described in~\ref{sec:method:architecture:baselines}.

\subsubsection{Hyperparameter Tuning}
\label{sec:method:experiment:hp}
Hyperparameters were selected by an Optuna \citep{akiba2019optuna} sweep on the held-out
balanced subset of~\ref{sec:method:data:hpsubset}, excluded from all subsequent training and
evaluation so reported scores cannot absorb information from tuning, at a cost of 50,000
graphs. Tuning was performed once, on the full-graph configuration, with values held fixed
for every reduced configuration: re-tuning per method would confound retained signal with
tuning effort. Reduced configurations are therefore evaluated at hyperparameters chosen for
a different input distribution, which if anything biases against them.

% Original, condensed above:
% An Edge-Aware Graph Convolutional Network (GCN+) is trained to predict the optimizability of
% a circuit when subjected to the \texttt{Orchestrate} synthesis script. The architecture is
% described in~\ref{sec:method:architecture} and is held fixed across all reduction
% configurations; the baselines it is compared against are described
% in~\ref{sec:method:architecture:baselines}.
%
% Hyperparameters were selected by an Optuna \citep{akiba2019optuna} sweep on the held-out
% balanced subset of~\ref{sec:method:data:hpsubset}, which is then excluded from all
% subsequent training and evaluation. Tuning on a subset that later appeared in training or
% test would let the reported scores absorb information from the tuning process; excluding it
% entirely costs 50,000 graphs and removes the question.
%
% Tuning was performed once, on the full-graph configuration, and the resulting values are held
% fixed for every reduced configuration. This is the conservative choice for the comparison
% being made: re-tuning per reduction method would confound ``this reduction retains more
% signal'' with ``this reduction received more tuning effort''. It does mean reduced
% configurations are evaluated at hyperparameters chosen for a different input distribution,
% which if anything biases against them, and bounds the strength of any negative result about
% reduction.

The search covered architecture and optimization jointly, 13 to 18 dimensions depending on
the encoder drawn. Architectural dimensions: encoder (six candidates, a graph convolutional
network, two edge-conditioned isomorphism networks, a plain message-passing network, an
attentional convolution and a hybrid local/global model), hidden width in
$\{32, 64, 128, 256, 512\}$, depth $L \in \{2, \dots, 8\}$, jumping-knowledge mode in
$\{\text{last}, \text{cat}, \text{max}, \text{sum}\}$, normalisation in
$\{\text{batch}, \text{layer}, \text{graph}\}$, pooling in $\{\text{mean}, \text{max},
\text{sum}\}$, positional encoding in $\{\text{none}, \text{level}, \text{input paths},
\text{local shortest-path sum}\}$ with width in $\{16, 32, 64, 128\}$, and, where applicable,
attention heads in $\{1, 2, 4\}$, edge-channel width in $\{32, 64, 128\}$ and update-layer
depth in $\{2, 3, 4\}$. Optimization dimensions: learning rate and weight decay (each over
about two decades, $10^{-4}$ to $6 \times 10^{-3}$ and $3 \times 10^{-5}$ to $10^{-2}$),
dropout, the loss transition point, and batch size in graphs.
% TODO: the sweep searched the loss transition point; the reported value was set by hand.

The sweep ran in two stages across three parallel workers with distinct sampler seeds and a
shared data split: an exploration stage requesting 50 trials per worker on a 15,000-graph
subsample (one-hour cap), and an exploitation stage requesting 20 per worker on 35,000
graphs (four-hour cap), warm-started from the ten best exploration configurations. Of the
210 trials requested, the surviving worker logs record 94 started, 52 completed and 38
pruned, 34 of those on memory exhaustion: the tuning stage hits the same bottleneck the
thesis is about, pruning hardest on precisely the configurations that would have been most
expensive to train.

% Original, condensed above:
% The search covered architecture and optimization jointly, in 13 to 18 dimensions depending
% on which encoder a trial drew. The architectural dimensions were the encoder itself, drawn
% from six candidates (a graph convolutional network, two edge-conditioned isomorphism
% networks, a plain message-passing network, an attentional convolution and a hybrid
% local/global model), hidden width in $\{32, 64, 128, 256, 512\}$, message-passing depth
% $L \in \{2, \dots, 8\}$, jumping-knowledge mode in $\{\text{last}, \text{cat}, \text{max},
% \text{sum}\}$, normalisation in $\{\text{batch}, \text{layer}, \text{graph}\}$, graph
% pooling in $\{\text{mean}, \text{max}, \text{sum}\}$, positional encoding in
% $\{\text{none}, \text{level}, \text{input paths}, \text{local shortest-path sum}\}$ with
% width in $\{16, 32, 64, 128\}$, and, on the encoders that have them, attention heads in
% $\{1, 2, 4\}$, edge-channel width in $\{32, 64, 128\}$ and update-layer depth in
% $\{2, 3, 4\}$. The optimization dimensions were learning rate, weight decay, dropout, the
% loss transition point and the batch size in graphs, the first two searched over about two
% decades each ($10^{-4}$ to $6 \times 10^{-3}$ and $3 \times 10^{-5}$ to $10^{-2}$).
%
% The sweep ran in two stages across three parallel workers with distinct sampler seeds and a
% shared data split. An exploration stage requested 50 trials per worker on a 15,000-graph
% subsample with a one-hour cap per trial; an exploitation stage requested 20 per worker on
% 35,000 graphs with a four-hour cap, warm-started from the ten best exploration
% configurations. Of the 210 trials requested, the surviving worker logs record 94 started, 52
% completed and 38 pruned, 34 of those on memory exhaustion. The tuning stage therefore hits
% the same bottleneck the thesis is about, and it prunes hardest on precisely the
% configurations that would have been most expensive to train.

\subsubsection{Training Configuration}
\label{sec:method:experiment:training}
Smooth L1 loss, transition point $\beta = 0.01$~\eqref{eq:method:smoothl1}, set well below
the default since the target occupies a narrow band within $[0, 1]$. Optimization uses AdamW
at learning rate $3 \times 10^{-4}$, weight decay $3.96 \times 10^{-3}$, linear warmup over
10,000 steps and plateau-triggered decay thereafter, stopping early on validation loss
(remaining schedule, clipping and regularisation values in~\ref{sec:apx:config}). Two
settings change the measurements, not the model: the training step is compiled (moving
every timing figure) and mixed precision is used (moving every memory figure), disclosed
here since no absolute number means anything without them.

% Original, condensed above:
% Smooth L1 loss with transition point $\beta = 0.01$~\eqref{eq:method:smoothl1}, set well
% below the default because the target occupies a narrow band within $[0, 1]$ and the default
% would place nearly every residual in the quadratic branch. Optimization uses AdamW at
% learning rate $3 \times 10^{-4}$ and weight decay $3.96 \times 10^{-3}$, with linear warmup
% over 10,000 steps and plateau-triggered decay thereafter, stopping early on validation loss.
% The remaining schedule, clipping and regularisation values are listed
% in~\ref{sec:apx:config}.
%
% Two settings change the measurements rather than the model and are disclosed with them: the
% training step is compiled, which moves every timing figure, and mixed precision is used on
% the accelerator, which moves every memory figure. Neither varies across configurations, so
% the comparisons hold, but no absolute number here means anything without them.
% Optimization uses AdamW with learning rate $3 \times 10^{-4}$ and weight decay
% $3.96 \times 10^{-3}$, a linear warmup over the first 10,000 steps, and
% \texttt{ReduceLROnPlateau} thereafter (factor 0.5, patience 2, floor $10^{-6}$). Training
% stops early on validation loss with patience 3. Gradients are clipped at 1.0.

\subsubsection{Execution Environment and Workload Profile}
\label{sec:method:experiment:hardware}
All experiments ran on the Snellius national supercomputer: H100 80\,GB GPU nodes for
training and evaluation, CPU-only nodes for the offline reduction precompute (embarrassingly
parallel, no GPU benefit). Software versions and wall-clock budget are recorded
in~\ref{sec:method:experiment:reproducibility}.

Two measured workload characteristics explain the shape of the RQ2 results. Memory scales
close to linearly with nodes and edges per batch, the encoder's own parameters negligible at
this width and depth, which is why a node budget bounds memory effectively. Message passing
is latency-bound on gather/scatter operations rather than compute-bound, floating-point
pipelines and bandwidth both sitting far below saturation at full occupancy, so a larger
batch amortises overhead without a proportional speedup, and reductions that remove
\emph{edges} should help throughput more than their node compression alone suggests.

% Original, condensed above:
% All experiments ran on the Snellius national supercomputer: H100 80\,GB GPU nodes for
% training and evaluation, and CPU-only nodes for the offline reduction precompute, which is
% embarrassingly parallel over graphs and does not benefit from a GPU\@. Software versions and
% the total wall-clock budget are recorded in~\ref{sec:method:experiment:reproducibility}.
%
% Two measured characteristics of the workload belong here rather than there, because they
% explain the shape of the RQ2 results. Memory scales close to linearly with nodes and edges
% per batch, the encoder's own parameters being negligible at this width and depth, which is
% why a node budget bounds memory effectively at all. Message passing is latency-bound on
% gather and scatter operations rather than compute-bound: at full device occupancy the
% floating-point pipelines and memory bandwidth both sit far below saturation. A larger batch
% therefore amortises per-batch overhead without delivering a proportional speedup, and
% reductions that remove \emph{edges} should help throughput more than their node compression
% alone would suggest.